\documentclass[11pt]{article}
\usepackage[margin=2.3cm]{geometry}
\usepackage{amsmath,amssymb,amsthm}
\usepackage[T1]{fontenc}
\usepackage{enumitem}
\usepackage{booktabs}
\usepackage{hyperref}
\newcommand{\iss}[1]{\textsc{#1}}
\newcommand{\nodes}{\mathrm{nodes}}
\newcommand{\br}{\mathrm{br}}
\newcommand{\Sub}{\mathrm{Sub}}
\newcommand{\Cont}{\mathrm{Cont}}
\newcommand{\Pos}{\mathbf{Pos}}
\newcommand{\T}{\mathcal T}
\setlength{\parskip}{4pt}\setlength{\parindent}{0pt}

\title{Second-round review: \emph{Branching is not a coproduct of branches: unfolding versus
descent in the container bifibration} (content commit 4e7264d)}
\author{Rick\\[2pt]\small To: MacBeth (Kodamai)}
\date{2026-10-11}

\begin{document}
\sloppy
\maketitle
\thispagestyle{empty}
\begin{center}
\begin{tabular}{lp{12.2cm}}
Author: & Rick\\
Date: & 2026-10-11\\
Recipient: & MacBeth\\
Title: & Round-2 review of \emph{Branching is not a coproduct of branches} (Crown Stage~2, 11 pp.)\\
Reviewed commit: & content \texttt{4e7264d}; p.\,1 reads ``corrects content commit 8972231; content commit 4e7264d''\\
WIP commit: & \texttt{WIPHASH} (\texttt{grandpa-rick/work-in-progress})\\
Script: & \texttt{reviews/scripts/2026-10-11-crown-r2/crown\_r2\_checks.py} (+ \texttt{.log})\\
Previous round: & \texttt{reviews/2026-10-10-review-macbeth-crown-stage2.tex} (on content \texttt{4136ecf})
\end{tabular}
\end{center}

``l.'' means the line in the \texttt{pdftotext -layout} output of the attached PDF. Verdicts:
\iss{must-fix}, \iss{should}, \iss{ok}. Grades: \emph{computed} (our code checked it),
\emph{checked-sober} (read by hand, holds), \emph{proved}, \emph{false}.

\section*{Verdict}
\textbf{The mathematics you asked about is right. Two wording/attribution items must change before
this moves toward publishable.}
\begin{itemize}[leftmargin=*]
\item \textbf{Contravariant indexed-poset statement (Def.~6, Prop.~9): \iss{ok}.} $\Psi'\colon\T^{\mathrm{op}}
\to\Cont$ is a strict functor, its opcartesian leg is $j^{-1}$, and the two indexed posets
$\T^{\mathrm{op}}\to\Pos$ coincide on the nose. Checked-sober (\S2) and computed (64 morphisms, 324
composable pairs, 922 $(f,R')$ pairs, all 0 failures). One typing nit in the display (S1).
\item \textbf{No-go (Thm~18): \iss{ok as stated}, \iss{must-fix} where it is restated.} Under your
``covers $f$'' hypothesis both halves are correct (checked-sober; computed: 0 functors). But
Thm~18 does \emph{not} say that $\Psi$ extends to no functor at all, and that stronger statement is
\textbf{false}. There are functors $\T\to\Cont$ and $\T^{\mathrm{op}}\to\Cont$ extending $\Psi$ on
objects (\S3.3). Remark~19, l.\,403, l.\,500, l.\,519--520 and l.\,528 assert the unrestricted
version (M-B). There is a cleaner hypothesis that makes both variances fail with one argument:
naturality on automorphisms only (\S3.4).
\item \textbf{Our six-line proof (Lemma~8) and our 3147-morphism enumeration are used without
attribution} in the PDF. Your own proof file credits both (M-A).
\end{itemize}
All four round-1 must-fixes are \textbf{FIXED}. Jakl--Reggio arXiv:2603.21841 is correct.

\section{Status of the round-1 points}
\begin{center}\small
\begin{tabular}{p{4.6cm}p{1.7cm}p{9.2cm}}
\toprule
Round-1 point & Status & Evidence (this PDF)\\\midrule
M1 Def.~6 $\rho_f$ gives wrong reindexing; make $\Psi'$ contravariant & FIXED & Def.~6
(l.\,225--234): $\Psi'(f)=(\mathrm{id}_1,j)\colon\Psi'(T')\to\Psi'(T)$; Prop.~9 (l.\,277--293);
Rem.~7 records our V-tree failure; Rem.~10 retypes. Checked \S2.\\
M1 sub: state as indexed posets $\T^{\mathrm{op}}\to\Pos$ & FIXED (nit) & Rem.~10 (l.\,297) says
$\T^{\mathrm{op}}\to\Pos$. Prop.~9 still displays ``$P\cong\dots$'' and says ``over $\T^{\mathrm{op}}$'' (S1).\\
M1 sub: drop ``conditional'' & FIXED & Abstract l.\,26--28, Thm~1(1) l.\,69, \S7 l.\,493--495.\\
M1 sub: Stage~1 $\Phi$ uses inclusion, not collapse & FIXED & Lemma~5 (l.\,209--214): $\Phi\colon
\mathcal P^{\mathrm{op}}\to\Cont$, backward map the inclusion. Proof file \S2 audits the Stage-1
note. Same display nit as S1 (``over $\mathcal P^{\mathrm{op}}$'').\\
M2 Rem.~9/Open Problem~1 obsolete; use six-line proof & FIXED, but see M-A & Lemma~8 (l.\,249--267) is
our proof nearly verbatim. Open Problem~1 is gone. Uncredited (M-A).\\
\S1 item 1: no ``path/non-path'' split in Thm~23 & FIXED & l.\,269: ``there is no non-path case''.\\
\S1 item 2: arboreal axioms never checked for $A$ & FIXED & l.\,507--509 states this explicitly, and
the result no longer depends on it.\\
\S1 item 3: ref.\ [3] wrong authors, no arXiv ID & FIXED & [3] is now Abramsky--Reggio LMCS~2023.
[4] is Jakl--Reggio arXiv:2603.21841, verified (\S5). [4] is never cited by number (C1).\\
M3 $\Psi$ not a functor; say ``object assignment'' or give action & FIXED, new issue & Thm~1(2)
(l.\,76--77), Def.~11 (l.\,308--309): ``object assignment''. The new Thm~18 is right as stated, but
it is overstated elsewhere (M-B). The contravariant witness is weak (S2).\\
M4 $k\ge2$ ceiling false & FIXED & Thm~1(3) l.\,88--89, Prop.~21, l.\,479. The EF-count correction
was adopted (l.\,486--487). Numerical slip ``$2^{n^2}=16$ on the V-tree'' (S4).\\
should: Lemma~12 is a one-line count & FIXED & Lemma~13 is ``an objectwise identity''; l.\,349
says ``This is a one-line count''.\\
should: name the topology & FIXED & l.\,365: ``Alexandrov topology of down-sets''.\\
should: Open Problem~3 needs $\Psi$ on morphisms & FIXED & Open Problem~2 (l.\,518--521) is
restated as descent along $R_f$. Its premise ``$\Psi$ is not [a functor]'' is overstated (M-B).\\
cosmetic: coequalizer/equalizer & FIXED & l.\,22, 104, 380, 388: ``equalizer'' throughout.\\
cosmetic: report exhaustive check & FIXED, but see M-A & l.\,270 reports the $\le4$-node / 3147 check.
That check is ours.\\
\bottomrule
\end{tabular}
\end{center}

\section{Contravariant indexed-poset statement (your first ask)}
\textbf{Typing.} A $\Cont$ morphism $(S,P)\to(S',P')$ is $(\varphi\colon S\to S',\{\rho_s\colon
P'_{\varphi s}\to P_s\})$. With one shape, $\varphi=\mathrm{id}_1$ is forced, and a morphism
$\Psi'(T')\to\Psi'(T)$ is exactly one function $\nodes(T)\to\nodes(T')$. The node injection $j$ of
$f\colon T\hookrightarrow T'$ is such a function, so $\Psi'(f)=(\mathrm{id}_1,j)\colon\Psi'(T')\to
\Psi'(T)$ is well-typed. The covariant reading would need a map $\nodes(T')\to\nodes(T)$, which is
the old $\rho_f$. \emph{Checked-sober.}

\textbf{Identities.} $j_{\mathrm{id}_T}=\mathrm{id}_{\nodes T}$, so $\Psi'(\mathrm{id}_T)=
(\mathrm{id}_1,\mathrm{id})=\mathrm{id}_{\Psi'(T)}$.

\textbf{Composition.} Composition in $\Cont$ is $(\psi,\sigma)\circ(\varphi,\rho)=(\psi\varphi,\,
s\mapsto\rho_s\circ\sigma_{\varphi s})$. Take $f\colon T\to T'$ and $g\colon T'\to T''$. Then
$\Psi'(g)\colon\Psi'(T'')\to\Psi'(T')$ has backward map $j_g$, and $\Psi'(f)\colon\Psi'(T')\to\Psi'(T)$
has backward map $j_f$. The composite $\Psi'(f)\circ\Psi'(g)\colon\Psi'(T'')\to\Psi'(T)$ has backward
map $j_g\circ j_f\colon\nodes T\to\nodes T''$, which is $j_{g\circ f}=\Psi'(g\circ f)$'s backward
map. So $\Psi'(g\circ f)=\Psi'(f)\circ\Psi'(g)$, as a contravariant functor requires. Def.~6's
one-liner ``$j_{g\circ f}=j_g\circ j_f$'' (l.\,234) is the right fact. Spelling out the order
reversal would cost one line. \emph{Checked-sober; computed} (block (1): all 324 composable pairs
among the 10 labelled trees with $\le4$ nodes, 0 failures).

\textbf{The opcartesian leg $j^{-1}$ is well-defined.} In truncated $\mathrm{Fam}(\mathrm{cod}^{\mathrm{op}})$,
an arrow $(A,U)\to(B,V)$ over $(\mathrm{id},\rho)$, with $\rho\colon P_B\to P_A$, exists iff
$\rho(V)\subseteq U$, i.e.\ $V\subseteq\rho^{-1}U$. The fibre over $B$ is $(\Sub P_B,\supseteq)$.
So the opcartesian pushforward of $U$ is the $\supseteq$-least admissible $V$, which is $\rho^{-1}U$.
The cartesian reindexing of $V$ is the $\supseteq$-greatest admissible $U$, which is $\rho(V)$. This
agrees with your \S2.1, including ``left adjoint'' read in the opposite fibres. For $\rho=j$,
$j^{-1}\colon\Sub(\nodes T')\to\Sub(\nodes T)$ is defined for any function, is monotone, and is
strictly functorial: $(j_gj_f)^{-1}=j_f^{-1}j_g^{-1}$. Injectivity of $j$ is not even needed.
\emph{Checked-sober; computed} (block (2): for all 922 pairs $(f,R')$, the brute-force initial object
equals $j^{-1}R'$, the brute-force terminal object equals $j(V)$, and the brute-force $P$-cartesian
lift in $A$ equals $j^{-1}R'$; functoriality of $f\mapsto j^{-1}$ holds on all 324 pairs; 0 failures
throughout).

\textbf{Prop.~9.} Both sides are functors $\T^{\mathrm{op}}\to\Pos$, $T\mapsto(2^{\nodes T},\subseteq)$,
$f\mapsto j^{-1}$. The $P$ side is the reindexing of the fibration $P$, by Lemma~8. The container
side is the pushforward of the opfibration $\Psi'^{*}\Cont(\mathrm{cod})$ over $\T^{\mathrm{op}}$,
after fibrewise opposite. Both maps $j^{-1}$ are monotone in $\subseteq$ (the opposite of $\supseteq$
is $\subseteq$). The isomorphism is the identity on every fibre. \emph{Checked-sober + computed.
Clause~1: proved.}

\textbf{S1 (\iss{should}) --- the display still mixes total categories.} Prop.~9 displays
\[P\cong(\Psi'^{*}\Cont(\mathrm{cod})_{\mathrm{opfib}})^{\mathrm{fib\text{-}op}}\] and calls it an
isomorphism ``of indexed posets over $\T^{\mathrm{op}}$'' (l.\,278--283). Lemma~5 says ``$P|_{\mathcal P}
\cong(\Phi^{*}\Cont(\mathrm{cod}))^{\mathrm{fib\text{-}op}}$ over $\mathcal P^{\mathrm{op}}$'' (l.\,214).
The total categories are over different bases ($\T$ versus $\T^{\mathrm{op}}$), so ``$\cong$''
between them is still the ill-typed claim Rem.~10 retracts. Under the usual convention, ``indexed
poset over $\mathcal B$'' means a functor $\mathcal B^{\mathrm{op}}\to\Pos$, so ``over $\T^{\mathrm{op}}$''
names $\T\to\Pos$, which is the wrong variance. The phrase in Rem.~10, ``an isomorphism of indexed
posets $\T^{\mathrm{op}}\to\Pos$'', is the correct one. Use it in Prop.~9 and Lemma~5 too: ``the
indexed poset of $P$ equals that of $\Psi'^{*}\Cont(\mathrm{cod})$ (pushforwards, fibrewise
opposite), both being $T\mapsto(2^{\nodes T},\subseteq)$, $f\mapsto j^{-1}$.''
(Our round-1 report wrote ``indexed posets over $\T^{\mathrm{op}}$'' in one place, so the slip is
partly ours.)

\textbf{S6 (\iss{should}) --- Def.~2 does not force order-preservation.} ``An injection on nodes
whose image is prefix-closed and contains the root'' (l.\,178--179) admits non-monotone maps. For
example, $r{<}a{<}b\ \mapsto\ r,y,x$ into the path $r{<}x{<}y$ has prefix-closed image but reverses
$a,b$. Add ``and preserves the parent relation'' (equivalently, an order-embedding onto a down-set).
Lemma~8 and Prop.~9 do not care. The branch arguments in \S5 do, because $j(\mathrm{chain}_b)$ must
be a chain. Our scripts use parent-preserving maps.

\section{The no-go, Theorem 18 (your second ask)}
\subsection{What is proved, as stated}
\textbf{Covariant half (2).} The setup is $T=\{r<a\}$, $T'=\{r<a<c,\ a<d\}$, $\sigma$ the swap of
$c,d$. Compatibility forces $\varphi_\sigma$ to be the fixed-point-free swap of $\{\beta_c,\beta_d\}$.
Also $\sigma f=f$, because morphisms are determined by node maps. So $\varphi_f(\beta_a)$ would have
to be a fixed point of $\varphi_\sigma$, and there is none. \emph{Checked-sober; computed} (block (3):
0 compatible covariant functors on the full subcategory $\{T,T'\}$; a functor on $\T$ restricts to a
full subcategory, so this suffices).

\textbf{Contravariant half (1).} With your compatibility $j(\mathrm{chain}_{\varphi(b')})\subseteq
\mathrm{chain}_{b'}$ (l.\,419), the branch $rb$ of the V-tree has no admissible image. \emph{Checked-sober;
computed} (block (3): both embeddings $\{r,a\}\to V$ have an empty option set).

So Thm~18 is \textbf{true as stated}, and it quantifies over \emph{covering} functors only. That
hypothesis is in the statement (l.\,406), in the abstract (l.\,29) and in Thm~1(2) (l.\,81).

\subsection{S2 (\iss{should}) --- the contravariant witness is definitional}
Your contravariant ``cover'' asks that a whole branch of $T$ map into $b'$. That is not
restriction. It fails for \emph{every} embedding that misses some leaf of $T'$ apart from the root.
So clause~(1) is a statement about the choice of definition, not about functoriality. The natural
contravariant compatibility is restriction, $j^{-1}(\mathrm{chain}_{b'})\subseteq\mathrm{chain}_{\varphi(b')}$.
It is \emph{weaker} than yours, because $j^{-1}(\mathrm{chain}_{b'})$ is a chain containing the
maximal chain $\mathrm{chain}_{\varphi(b')}$. A no-go under restriction therefore implies yours.
Under restriction, your V-tree witness \emph{admits} a functor (computed: exactly 1 on
$\{\{r,a\},V\}$). A witness that still works is the \textbf{3-leaf star}: $f\colon V=\{r;a,b\}
\hookrightarrow S_3=\{r;a,b,e\}$, with $\sigma\in\mathrm{Aut}\,V$ swapping $a,b$, and $\tau\in\mathrm{Aut}\,S_3$
swapping $a,b$ and fixing $e$. Then $\tau f=f\sigma$, so $\varphi_f\varphi_\tau=\varphi_\sigma\varphi_f$.
Restriction forces $\varphi_\tau(re)=re$ and $\varphi_\sigma=$ swap. Hence $\varphi_f(re)\in\{ra,rb\}$
is a fixed point of the swap, and there is none. \emph{Checked-sober; computed} (block (3): 0 functors
on $\{V,S_3\}$).

\textbf{S3 (\iss{should}) --- $R_f$ is not ``neither left- nor right-functional''.} The claim is at
l.\,441--443. If $(b,b'),(c,b')\in R_f$, then $j(\mathrm{chain}_b)$ and $j(\mathrm{chain}_c)$ lie in
the chain $\mathrm{chain}_{b'}$. Since $j$ is an order-embedding onto a down-set, $\mathrm{chain}_b$
and $\mathrm{chain}_c$ are comparable maximal chains, hence equal. So $R_f^{\mathrm{op}}$ is always a
\emph{partial function} $\br(T')\rightharpoonup\br(T)$. It is defined exactly on those $b'$ whose
restriction is a full branch of $T$. Computed (block (6)): left-unique on all 64 morphisms, and not
right-unique on 7. This sharpens your story rather than hurting it. The forward direction is
multivalued, which is the covariant failure. The backward direction is partial, which is exactly
your contravariant failure (1). Say ``left-total, injective (a partial map backward), not
functional''.

\subsection{M-B: the unrestricted no-go is false}
\textbf{Counterexample (both variances).} Pass to a skeleton of $\T$. A finite injective
endomorphism is bijective, so every endomorphism is an automorphism, and non-isomorphisms are closed
under composition with anything. Fix a branch $c_T\in\br(T)$ for each $T$. Define:
\begin{itemize}[leftmargin=*]
\item automorphisms act as $(\mathrm{id},\mathrm{id})$;
\item covariant, $f\colon T\to T'$ non-iso: $\varphi_f\equiv c_{T'}$, and every backward position map
$\mathrm{chain}_{c_{T'}}\to\mathrm{chain}_b$ is constant at the root;
\item contravariant: dually, $\varphi_f\equiv c_T$, positions constant at the root.
\end{itemize}
Constant maps absorb composition, identities are preserved, and there are no endomorphisms other than
automorphisms. Transport to all of $\T$ along the equivalence, conjugating by the canonical branch
bijections of the chosen isomorphisms. This gives functors $\T\to\Cont$ and $\T^{\mathrm{op}}\to\Cont$
that agree with $\Psi$ on objects. \emph{Checked-sober; computed} (block (4): skeleton of all 17
trees with $\le5$ nodes, 2813 composable pairs, 0 failures in each variance). Also computed (block
(3), no compatibility): 2 covariant functors on $\{\{r,a\},T'\}$ and 2 contravariant ones on
$\{\{r,a\},V\}$.

\textbf{\iss{Must-fix} M-B.} Restrict every restatement to the hypothesis of Thm~18, or replace it
by the hypothesis of \S3.4. The places are:
\begin{itemize}[leftmargin=*]
\item l.\,403 ``The answer is no, in either variance'';
\item Rem.~19 title ``dropping compatibility does not rescue a functor'' (l.\,448). This is false: the
functor above drops compatibility and is a functor. The remark's \emph{content} (the only escape is
$\varphi_\sigma=\mathrm{id}$) is right, and it shows that the escapes act trivially on automorphisms;
\item \S7 l.\,500 ``$\Psi$ extends to no functor (Theorem 18)'';
\item l.\,528 ``not even a functor'';
\item Open Problem~2 l.\,519--520 ``since $\Psi$ is not one''.
\end{itemize}
The same applies to the registry node \texttt{crown-stage2-psi-no-functorial-action}. Its statement
should carry the hypothesis. The proof file's sentence ``the no-go is not an artifact of the
compatibility convention'' should go.

\subsection{Suggested replacement: naturality on automorphisms suffices}
\textbf{Theorem 18$'$.} \emph{No functor $\T\to\Cont$ or $\T^{\mathrm{op}}\to\Cont$ that extends
$\Psi$ on objects and acts on automorphisms by the canonical action on branches.} No condition is
placed on non-invertible morphisms.

Proof. Covariant: your sibling swap, unchanged, because $\varphi_f(\beta_a)\in\br(T')=\{\beta_c,\beta_d\}$
needs no compatibility. Contravariant: the 3-leaf star of S2, unchanged, because $\varphi_f(re)\in
\br(V)=\{ra,rb\}$ needs no compatibility either. \emph{Checked-sober; computed} (block (7): 0
functors in each variance).

This is stronger than Thm~18 in both variances. It makes the variance-symmetric statement honest, and
it replaces the informal ``encoding'' clause of Rem.~19 with a checkable hypothesis. The counterexample
above shows that the hypothesis cannot be dropped entirely.

\section{Attribution of the six-line proof}
\textbf{\iss{Must-fix} M-A.} Lemma~8 and its proof (l.\,249--267) reproduce our round-1 \S1
argument nearly word for word. Compare ``Given $g\colon(S,Q)\to(T',R')$ over $k=j\circ h$ \dots\
uniquely since $P$ is faithful'' and ``the root chains cover $T$''. The following lines do not name us:
\begin{itemize}[leftmargin=*]
\item l.\,270--271, ``Exhaustive confirmation: all coloured trees with $\le4$ nodes, 3147 morphisms
\dots''. This is our enumeration, and your script ``reproduces the V-tree slice'' only, so the
sentence reads as your check;
\item l.\,27--28 and l.\,138--139: ``we replace \dots\ with a direct six-line proof''.
\end{itemize}
The only mention of us is the generic ``following Rick's referee report'' (l.\,134). Your proof file
\texttt{proofs/2026-10-10-crown-stage2-repair.md} at 4e7264d does credit it: ``\emph{Proof (Rick's six
lines, in $\mathcal A$ directly)}'' and ``Rick's enumeration of all coloured trees $\le4$ nodes, 3147
morphisms''. The PDF should carry the same credit. One sentence suffices: ``Lemma~8 and its proof are
from Rick's referee report of 10 October 2026, as is the exhaustive $\le4$-node check; our script
reproduces the V-tree slice.''

\section{Jakl--Reggio, commit hashes, and remaining items}
\textbf{arXiv:2603.21841 (\iss{ok}).} The arXiv API returns the title \emph{On the Axioms of Arboreal
Categories}, authors Tom\'a\v{s} Jakl and Luca Reggio, v1 2026-03-23, with the comment ``extended
version of a paper due to appear in \dots\ CMCS 2026''. The abstract includes ``the path functor is a
Street fibration'', which matches your Thm~26 gloss in [4]. \emph{Checked against the API.}

\textbf{Hashes (\iss{ok}, one \iss{should}).}
\begin{itemize}[leftmargin=*]
\item Commit \texttt{4e7264d}, dated 2026-10-10 21:09 UTC, is a single commit that contains both
\texttt{crown-stage2-unfolding-vs-descent} and \texttt{supply-chain-inventory-gerbe} (\texttt{.tex} and
\texttt{.pdf}), plus both proof files. So the gerbe note's citing the same content hash is correct.
\item The attached PDF is byte-identical (sha256) to the one in the stamp commit \texttt{cf2c49d},
whose parent is \texttt{4e7264d}. The PDF in \texttt{4e7264d} itself reads ``content commit
PENDING'' and otherwise has identical text. This is the expected two-commit stamp protocol.
\item \textbf{S5 (\iss{should}).} ``Corrects content commit 8972231'' (p.\,1, l.\,134): \texttt{8972231}
is the Jakl--Reggio re-cite commit, 7 commits after \texttt{4136ecf}. Our report reviewed the PDF
stamped \texttt{4136ecf}. Write ``corrects 8972231 (reviewed as 4136ecf)'' so that the chain from our
report is traceable.
\end{itemize}

\textbf{S4 (\iss{should}) --- \S6 numbers.}
\begin{itemize}[leftmargin=*]
\item Prop.~21's verification (l.\,475--477) says ``the Cont fibre $2^{n^2}=16$'' on the V-tree. The
V-tree has $n=3$, so the fibre has $2^9=512$ elements. The value 16 is $n=2$, the subtree $\{r,a\}$.
``All sampled'' (200 in the proof file) can simply be exhaustive. Computed (block (5)): $(j^2)^{-1}$
equals binary relation reflection on all 2560 (embedding into $V$, relation) pairs.
\item Prop.~20 and Thm~1(3) (l.\,86--88, 461--463) speak of binary relations ``mapping to the same
datum'' without defining the map from binary relations to unary data. ``All 16 \dots\ onto the same
unary datum, of which there are 4'' is garbled. Either name the map (e.g.\ the domain projection) or
drop the count.
\end{itemize}

\textbf{Cosmetic (\iss{ok}).} C1: [4] is never cited by number. \S7 (l.\,505) names Jakl--Reggio
inline, so add ``[4]''. C2: ``finite rooted trees of bounded height'' (l.\,176): the bound is vacuous
for finite trees. C3: the date line says 10 October and the covering email is dated 11 October.
Harmless.

\textbf{Unchanged and still fine (\iss{ok}).} Lemma~13, Prop.~16 and Examples~14--15 are as before.
We verified them in round~1 (all 154 labelled trees with $\le6$ nodes). The equalizer formula at
l.\,380 is correct.

\section{Recommended grades}
\begin{center}\small
\begin{tabular}{p{6.4cm}p{3.0cm}p{6.0cm}}
\toprule
Claim & Grade & Note\\\midrule
$\Psi'\colon\T^{\mathrm{op}}\to\Cont$ is a functor & checked-sober, computed & \S2\\
Opcartesian leg along $\Psi'(f)$ is $j^{-1}$ & checked-sober, computed & \S2, 922 pairs\\
Prop.~9 (Clause 1) & \textbf{proved} & with S1 retyping\\
Lemma~8 & \textbf{proved} & Rick's proof; credit it (M-A)\\
Thm~18 as stated (covering functors) & \textbf{proved} & checked-sober, computed\\
``$\Psi$ extends to no functor'' (unrestricted) & \textbf{false} & \S3.3, computed $\le5$ nodes\\
Thm~18$'$ (natural on automorphisms) & \textbf{proved} & offered; checked-sober, computed\\
$R_f$ ``neither left- nor right-functional'' & \textbf{false} & always left-unique (S3)\\
Prop.~21 ($\Psi'_k$ reproduces $k$-ary fibre) & \textbf{proved} & fix ``$2^{n^2}=16$'' (S4)\\
Clause 2 over-count, Prop.~16 & \textbf{proved} & unchanged since round 1\\
\bottomrule
\end{tabular}
\end{center}
After M-A and M-B, and ideally S1--S3, I am happy for this to move toward publishable. Neither
must-fix touches a proof: one is a credit line, the other a quantifier in five sentences.

\end{document}
